\chapter{Backup}


\section{Observable requirements: Resummation chapter}
As a concrete example, consider thrust in the soft-collinear limit.
Evaluating $\tau = 1 - T$ when an additional emission is soft and collinear
to leg $\ell$, one finds~\cite{Banfi:2004yd}:
\begin{equation}
	\tau\simeq\tau_\mathrm{sc}=\sum_{\ell=1,2}\sum_{i\in\mathcal{H}_\ell}\frac{k_{ti}^{(\ell)}}{Q}e^{-\eta^{(\ell)}}\,,
\end{equation}
which is of the form of Eq.~\eqref{eq:general-sc-V}, with parameters
\begin{equation}
    a = b_\ell = d_\ell = g_\ell(\phi) = 1\,, \quad \ell = 1, 2\,.
\end{equation}
This particularly simple parameter set gives rise to a correspondingly 
simple resummation structure.

\section{NLL Resummation: Resummation chapter}

To evaluate $\Sigma(\{p\},v)$, we split the exponent in
Eq.~\eqref{eq:cumu-factor-more} according to whether the observable value
for a single emission $k$ exceeds $v$:
\begin{equation}
\label{eq:exp-rad}
  \int[dk]\,M^2(k)
  = \underbrace{\int[dk]\,M^2(k)\,\Theta\!\left(V(\{\tilde{p}\},k)-v\right)}_{\displaystyle R(v)}+ \int[dk]\,M_{\mathrm{sc}}^2(k)\,\Theta\!\left(v-V_{\mathrm{sc}}(\{\tilde{p}\},k)\right)\,.
\end{equation}
In the second term, the full matrix element $M^2$ is replaced by its
soft-collinear limit $M_{\mathrm{sc}}^2$, and $V$ by $V_{\mathrm{sc}}$.
This is valid because the constraint
$\Theta(v-V_{\mathrm{sc}})$ restricts the emission to the
soft-collinear region, where $M^2\simeq M_{\mathrm{sc}}^2$, with the
remaining corrections contributing only at NNLL accuracy.\footnote{These
corrections are accounted for by the terms
$\delta\mathcal{F}_{\rm hc}$, $\delta\mathcal{F}_{\rm sc}$, and
$\Delta\mathcal{F}_{\rm wa}$ introduced in the NNLL extension in
Ch.~\ref{ch:hhZjet}.}

\section{Validation: N-jettiness subtraction}

Ref.~\cite{Gaunt:2015pea} presents a subtraction method for
performing NNLO QCD calculations for arbitrary processes, exploiting
the factorisation of the $N$-jettiness distribution in the limit
$\mathcal{T}_N\to0$. The method uses SCET to determine the IR
singular contributions to $N$-jet cross sections analytically and
constructs suitable $\mathcal{T}_N$-subtractions from them. The
singular structure is organised into hard, beam, jet, and soft
functions. The beam function $B_a$ encodes all collinear emissions
from incoming parton $a$ and depends on its flavour $\kappa_a$ and
momentum fraction $x_a$. The jet function $J_c$ encodes collinear
emissions from the outgoing parton $c$ and depends on its flavour
$\kappa_c$. The soft function $S_\kappa$, where
$\kappa\equiv\{\kappa_a,\kappa_b,\kappa_c\}$ denotes the full
partonic channel, describes all soft emissions between the partons.

Specialising this to the one-jettiness ($N=1$), we find that the singular part of the distribution ($\tau_1\to0$) factorises as 
\begin{equation}
	\label{eq:Cnm-def}
	\frac{\mathrm{d}\sigma^{\rm sing}(\{p\})}{\mathrm{d}\tau_1}
	=\sum_{m\geq0}\left[\mathcal{C}_{-1}^{(m)}(\{p\},M,\mu_R)\,\delta(\tau_1)
	+\sum_{n=0}^{2m-1}\mathcal{C}_n^{(m)}(\{p\},M,\mu_R)\,\mathcal{L}_n(\tau_1)\right]
	\left(\frac{\alpha_s(\mu_R)}{4\pi}\right)^m\,,
\end{equation}
where 
\begin{equation}
	\mathcal{L}_n(\tau_1)=\left[\frac{\theta(\tau_1)\ln^n(\tau_1)}{\tau_1}\right]_+\,.
\end{equation} 
are the standard plus distributions. 

The coefficients have the expansion
\begin{subequations}
\label{eq:subt-mn-exp}
\begin{align}
    &\mathcal{I}_{\kappa_1 k,n}^{(m)}(z,\lambda_F,\lambda_1) =
    \mathcal{I}_{\kappa_1 k,n}^{(m)}(z,\lambda_F)
    + \sum_{r=1}^{2m-1-n}\frac{(n+r)!}{n!\,r!}\,
    \mathcal{I}_{\kappa_1 k,n+r}^{(m)}(z,\lambda_F)\,\ln^r\!\lambda_1\,,\\
    &J_{\kappa_3,n}^{(m)}(\lambda_3) = J_{\kappa_3,n}^{(m)}
    + \sum_{r=1}^{2m-1-n}\frac{(n+r)!}{n!\,r!}\,J_{\kappa_3,n+r}^{(m)}\,\ln^r\!\lambda_3\,, \\
    &\widehat{S}_{\kappa,n}^{(m)}(\{\hat{q}_\ell\},\lambda_s) =
    \widehat{S}_{\kappa,n}^{(m)}(\{\hat{q}_\ell\})
    + \sum_{r=1}^{2m-1-n}\frac{(n+r)!}{n!\,r!}\,
    \widehat{S}_{\kappa,n+r}^{(m)}(\{\hat{q}_\ell\})\,\ln^r\!\lambda_s\,,
\end{align}
\end{subequations}
and analogously for $\mathcal{I}_{\kappa_2 k,n}^{(m)}(z,\lambda_F,\lambda_2)$. 

For the $\mathcal{C}_{-1}^{(m)}$ coefficients:
\begin{subequations}
\begin{align}
    &\mathcal{I}_{\kappa_1 k,-1}^{(m)}(z,\lambda_F,\lambda_1) =
    \mathcal{I}_{\kappa_1 k,-1}^{(m)}(z,\lambda_F)
    + \sum_{n=0}^{2m-1}\mathcal{I}_{\kappa_1 k,n}^{(m)}(z,\lambda_F)\,
    \frac{\ln^{n+1}\!\lambda_1}{n+1}\,,\\
    &J_{\kappa_3,-1}^{(m)}(\lambda_3) =J_{\kappa_3,-1}^{(m)}+\sum_{n=0}^{2m-1}J_{\kappa_3,n}^{(m)}\,\frac{\ln^{n+1}\!\lambda_3}{n+1}\,, \\
   & \widehat{S}_{\kappa,-1}^{(m)}(\lambda_s) =
    \widehat{S}_{\kappa,-1}^{(m)}+\sum_{n=0}^{2m-1}\widehat{S}_{\kappa,n}^{(m)}\,
    \frac{\ln^{n+1}\!\lambda_s}{n+1}\,,
\end{align}
\end{subequations}


\subsection{$h_{10}$}
These are complete form of the coefficients when everything has been substituted in and simplified by 
Mathematica while keeping the argument of the logarithms as dimensionless. 
\begin{equation}
\begin{aligned}
\mathcal{C}^{(1)}_{-1} \;=\;& \frac{2}{3}\left(5C_F-3C_A\right)\pi^2 \,-\,4\pi\beta_0\ln\frac{M\,Q_{ab}^2\,Q_c}{\mu_R^4}\,+\,\frac{40\pi}{3}\beta_0 \\[1.5ex]
&+\sum_{k}\int\frac{dz}{z}\,\frac{f_k(x_a/z,\mu_F)}{f_a(x_a,\mu_F)}\,2I_{qk}^{(1)}
 \;+\;\sum_{k}\int\frac{dz}{z}\,\frac{f_k(x_b/z,\mu_F)}{f_b(x_b,\mu_F)}\,2I_{\bar{q}k}^{(1)} \\[1ex]
&+\sum_{k}\int\frac{dz}{z}\left(2P_{qk}\ln\frac{Q_a\,M}{Q^2}+2P_{qk}\ln\frac{Q^2}{\mu_F^2}+3C_F\ln\frac{Q^2}{Q_a\,M}\right)\frac{f_k(x_a/z,\mu_F)}{f_q(x_a,\mu_F)} \\[1ex]
&+\sum_{k}\int\frac{dz}{z}\left(2P_{\bar{q}k}\ln\frac{Q_b\,M}{Q^2}+2P_{\bar{q}k}\ln\frac{Q^2}{\mu_F^2}+3C_F\ln\frac{Q^2}{Q_b\,M}\right)\frac{f_k(x_b/z,\mu_F)}{f_{\bar q}(x_b,\mu_F)} \\[1.5ex]
&+(C_A-2C_F)\left(\ln^2\frac{Q_{ab}^2}{Q_aQ_b}+2\left(I_{ab,c}+I_{ba,c}\right)\right) \\[1ex]
&-C_A\left(\ln^2\frac{Q_{ac}^2}{Q_aQ_c}+2\left(I_{ac,b}+I_{ca,b}\right)\right) \\[1ex]
&-C_A\left(\ln^2\frac{Q_{bc}^2}{Q_bQ_c}+2\left(I_{bc,a}+I_{cb,a}\right)\right) \\[1.5ex]
&+C_F\left[\ln^2\frac{Q_{ac}^2}{Q_{bc}^2}-\ln^2\frac{Q_{ab}^2}{\mu_R^2}
 +2\ln^2\frac{Q_a}{\mu_R}+2\ln^2\frac{Q_b}{\mu_R}
 +\ln^2\left(\frac{Q_aQ_b\,\mu_R^2}{Q_{ab}^4}\right)
 -\ln^2\left(\frac{M^2Q_aQ_b}{Q_{ab}^4}\right) \right.\\[1ex]
&\left.\hphantom{+C_F\left[\right.}
 -6\ln\frac{Q^2}{Q_{ab}\,\mu_R}-8
 +2\operatorname{Li}_2\left(\frac{-Q_{ac}^2}{M^2}\right)+2\operatorname{Li}_2\left(\frac{-Q_{bc}^2}{M^2}\right) \right.\\[1ex]
&\left.\hphantom{+C_F\left[\right.}
 +\ln^2\frac{M^2+Q_{ac}^2}{M^2}+\ln^2\frac{M^2+Q_{bc}^2}{M^2}
 -\ln^2\frac{Q_{ac}^2+M^2}{Q_{ac}^2}-\ln^2\frac{Q_{bc}^2+M^2}{Q_{bc}^2}\right] \\[1.5ex]
&+C_A\left[2\ln^2\left(\frac{\mu_R\,Q_{ab}^2\,Q_c}{Q_{ac}^2Q_{bc}^2}\right)+2\ln^2\frac{Q_c}{\mu_R}
 -\frac{1}{2}\ln^2\frac{Q_{ac}^2Q_{bc}^2}{Q_{ab}^2\mu_R^2}
 -2\ln^2\left(\frac{M\,Q_c\,Q_{ab}^2}{Q_{ac}^2Q_{bc}^2}\right)+\frac{11}{3} \right.\\[1ex]
&\left.\hphantom{+C_A\left[\right.}
 +2\operatorname{Li}_2\left(\frac{Q_{ab}^2-M^2}{Q_{ab}^2}\right)\right]
\end{aligned}
\end{equation}

The \textsc{ARES} coefficient $h_{10}$ for the $q\bar q\to g H$ channel is
\begin{equation}
\begin{aligned}
	h_{10}=\;&\frac{67\,C_A-10\,n_f}{18}
	+\frac{19\pi^2}{12}(2C_F-C_A)
	-2\pi\beta_0\ln\frac{M\,Q_c\,Q_{ab}^4}{\mu_R^6}\\[2ex]
	&+\sum_{k}\int\frac{dz}{z} \frac{f_k(x_a/z,\mu_F)}{f_a(x_a,\mu_F)}2I_{qk}^{(1)}+\sum_{k}\int\frac{dz}{z}\frac{f_k(x_b/z,\mu_F)}{f_b(x_b,\mu_F)}2I_{\bar{q}k}^{(1)}\\
	&+\sum_{k}\int\frac{dz}{z}\left(P_{qk}(z)\ln\frac{Q_a\,M}{Q^2}+P_{qk}(z)\ln\frac{Q^2}{\mu_F^2}+\frac{3}{2}C_F\ln\frac{Q^2}{Q_a\,M}\right)
    \frac{f_k(x_a/z,\mu_F)}{f_a(x_a,\mu_F)}\\
    &+\sum_{k}\int\frac{dz}{z}\left(P_{\bar{q}k}(z)\ln\frac{Q_b\,M}{Q^2}+P_{qk}(z)\ln\frac{Q^2}{\mu_F^2}+\frac{3}{2}C_F\ln\frac{Q^2}{Q_b\,M}\right)
    \frac{f_k(x_b/z,\mu_F)}{f_b(x_b,\mu_F)}\\[1ex]
	&+\frac{3}{2}C_F\ln\frac{Q_{ab}^2}{Q^2}
	+\frac{3}{2}C_F\ln\frac{Q_{ab}^2}{Q^2}\\[1ex]
	&+C_F\left[
	\ln^2\frac{M}{Q_a}
	+\ln^2\frac{M}{Q_b}
	-\ln^2\frac{Q_{ab}^2}{M^2}
	\right]
\\[1ex]
	&+\frac{C_A}{2}\left[
	2\ln^2\frac{M}{Q_c}
	+\ln^2\frac{M^2}{Q_{ab}^2}
	-\ln^2\frac{M^2}{Q_{ac}^2}
	-\ln^2\frac{M^2}{Q_{bc}^2}
	\right]\\[2ex]
	&+\frac{(C_A-2C_F)}{2}
	\left(
	\ln^2\frac{Q_{ab}^2}{Q_aQ_b}
	+2(I_{ab,c}+I_{ba,c})
	\right)\\[1ex]
	&-\frac{C_A}{2}
	\left(
	\ln^2\frac{Q_{ac}^2}{Q_aQ_c}
	+2(I_{ac,b}+I_{ca,b})
	\right)\\[1ex]
	&-\frac{C_A}{2}
	\left(
	\ln^2\frac{Q_{bc}^2}{Q_bQ_c}
	+2(I_{bc,a}+I_{cb,a})
	\right)\\
	&+\frac{C_A}{2}\left[
	\frac{152}{9}
	+4\operatorname{Li}_2\left(\frac{Q_{ab}^2-M^2}{Q_{ab}^2}\right)
	+2\ln\frac{Q_{ab}^2}{Q_{ac}^2}
	\ln\frac{Q_{ab}^2}{Q_{bc}^2}
	\right]\\[1ex]
	&+\frac{C_F}{2}\left(
	-16
	+4\operatorname{Li}_2\left(\frac{-Q_{ac}^2}{M^2}\right)
	+4\operatorname{Li}_2\left(\frac{-Q_{bc}^2}{M^2}\right)
	-2\ln^2\frac{Q_{ac}^2+M^2}{Q_{ac}^2}
	-2\ln^2\frac{Q_{bc}^2+M^2}{Q_{bc}^2}
\right.
\\[1ex]
	&\left.
	+2\ln^2\frac{M^2+Q_{ac}^2}{M^2}
	+2\ln^2\frac{M^2+Q_{bc}^2}{M^2}
	+2\ln^2\frac{Q_{ac}^2}{Q_{bc}^2}
	\right)
	-\frac{20}{18}n_f\,,
\end{aligned}
\end{equation}
where we have used the relation
\begin{equation}
\label{eq:I1-ARES}
I_{ik}^{(1)}(z) =
-\hat{P}_{ik}^{(\varepsilon)}(z)
+\hat{P}_{ik}(z;0)\!\left[
\left(\frac{\ln(1-z)}{1-z}\right)_{\!\!+}(1-z)
-\ln z\right]
+\frac{\pi^2}{6}\,\delta_{ij}\,C_i\,\delta(1-z)\,,
\end{equation}

\subsection{NNLO ingredients}
\label{subsec:NNLO-ingred}

Unlike the NLO case, there is no compact master expression for the
NNLO subtraction coefficients. For brevity the scale arguments are suppressed throughout this subsection. The contributions which make up these NNLO coefficients are organised into three types of quantity. The genuine $m$-loop contributions from the jet, beam, and soft functions as
\begin{equation}
\begin{aligned}
\label{eq:Xn-def}
    X_n^{(m)} \equiv &J_{c,n}^{(m)} + \widehat{S}_{\kappa,n}^{(m)} \\
&   + \sum_k\int\frac{dz}{z}\,\mathcal{I}_{ak,n}^{(m)}\frac{f_k(x_a/z,\mu_F)}{f_a(x_a,\mu_F)}
    + \sum_k\int\frac{dz}{z}\,\mathcal{I}_{bk,n}^{(m)}\frac{f_k(x_b/z,\mu_F)}{f_b(x_b,\mu_F)}\,.
    \end{aligned}
\end{equation}
The cross terms involving the one-loop
virtual amplitude $\vec{C}^{(1)}$ are
\begin{equation}
\label{eq:Xn11-def}
\begin{aligned}
    X_n^{(1+1)} \equiv& \,H^{(1)}_{\rm SCET}
    \left[J_{c,n}^{(1)} + \widehat{S}_{\kappa,n}^{(1)} \right.\\
  &\left.  + \sum_k\int\frac{dz}{z}\,\mathcal{I}_{ak,n}^{(1)}\frac{f_k(x_a/z,\mu_F)}{f_a(x_a,\mu_F)}
    + \sum_k\int\frac{dz}{z}\,\mathcal{I}_{bk,n}^{(1)}\frac{f_k(x_b/z,\mu_F)}{f_b(x_b,\mu_F)}\right]\,.
    \end{aligned}
\end{equation}
Finally, the cross-terms between two one-loop functions arising from the
$\mathcal{L}_n\otimes\mathcal{L}_m$ convolution are
\begin{equation}
\label{eq:Xnm11-def}
    X_{n,m}^{(1+1)} \equiv
J_{\kappa_3,n}^{(1)}\,\widehat{S}_{\kappa,m}^{(1)}+\frac{B_{\kappa_1}}{f_i}J_{\kappa_3,m}^{(1)}
    +\frac{B_{\kappa_2}}{f_j}J_{\kappa_3,m}^{(1)}+\frac{B_{\kappa_1}}{f_i}\widehat{S}_{\kappa,m}^{(1)}
    +\frac{B_{\kappa_2}}{f_j}\widehat{S}_{\kappa,m}^{(1)}+\frac{B_{\kappa_1}}{f_i}\frac{B_{\kappa_2}}{f_j}\end{equation}
With these definitions, the NNLO subtraction coefficients are
\begin{subequations}
\label{eq:NNLO-Cn}
\begin{align}
    \mathcal{C}_3^{(2)} &= X_3^{(2)} + X_{1,1}^{(1+1)}\,, 
    \label{eq:C32}\\
    \mathcal{C}_2^{(2)} &= X_2^{(2)} + \frac{3}{2}\!\left(
        X_{0,1}^{(1+1)} + X_{1,0}^{(1+1)}\right)\,,
    \label{eq:C22}\\
    \mathcal{C}_1^{(2)} &= X_1^{(2)} + X_1^{(1+1)} + 2X_{0,0}^{(1+1)}
        - \frac{\pi^2}{3}X_{1,1}^{(1+1)}+X_{1,-1}^{(1+1)}+X_{-1,1}^{(1+1)}\,,
    \label{eq:C12}\\
    \mathcal{C}_0^{(2)} &= X_0^{(2)} + X_0^{(1+1)}
        - \frac{\pi^2}{6}\!\left(X_{0,1}^{(1+1)}+X_{1,0}^{(1+1)}\right)
        + 2\zeta_3\,X_{1,1}^{(1+1)}+X_{0,-1}^{(1+1)}+X_{-1,0}^{(1+1)}\,.
    \label{eq:C02}
\end{align}
\end{subequations}
The last two terms Eqs.~\eqref{eq:C12} and~\eqref{eq:C02} are not present in Ref.~\cite{Gaunt:2015pea}. Please note that
we believe this to be an error on their side and these terms must contribute to make
up the full coefficient.

\subsection{NLO ingredients}
\label{subsec:NLO-ingred}


At NLO, the required coefficients are $\mathcal{C}_1^{(1)}$ and $\mathcal{C}_0^{(1)}$,
given by $n=1$ and $n=0$ respectively in
\begin{equation}
\label{eq:Cn1}
\begin{aligned}
    \mathcal{C}_n^{(1)}&(\{p\},M,\mu_R)=J_{c,n}^{(1)}(\lambda_c)
    +\widehat{S}^{(1)}_{\kappa,n}\!\left(\{\widehat{q}_\ell\},\lambda_s\right)\\
    &+\sum_{k}\int\frac{dz}{z}\,\mathcal{I}_{ak,n}^{(1)}(z,\lambda_F,\lambda_a)
    \frac{f_k(x_a/z,\mu_F)}{f_a(x_a,\mu_F)}
    +\sum_{k}\int\frac{dz}{z}\,\mathcal{I}_{bk,n}^{(1)}(z,\lambda_F,\lambda_b)
    \frac{f_k(x_b/z,\mu_F)}{f_b(x_b,\mu_F)}\,.
\end{aligned}
\end{equation}
The beam function coefficients are related to the
coefficients $\mathcal{I}_{ij}$ by
\begin{equation}
  B_{i,n}^{(m)}(x_i,\mu_R,\mu_F,\lambda_i) =
  \sum_k\int\frac{dz}{z}\,\mathcal{I}_{ik,n}^{(m)}\!\left(z,\lambda_F,\lambda_i\right)
  f_k\!\left(\frac{x_i}{z},\mu_F\right)\,.
\end{equation}
We have substituted these directly into Eq.~\eqref{eq:Cn1} for a
more transparent comparison with \textsc{ARES}, and refer to
$\mathcal{I}_{ij}$ as the beam function coefficients throughout.

The functions in Eq.~\eqref{eq:Cn1} are expanded in logarithms of the $\lambda$
parameters as
\begin{subequations}
\label{eq:subt-mn-exp}
\begin{align}
    J_{\ell=3,n}^{(m)}(\lambda_3) &= J_{\ell=3,n}^{(m)}
    + \sum_{r=1}^{2m-1-n}\frac{(n+r)!}{n!\,r!}\,J_{c,n+r}^{(m)}\,\ln^r\!\lambda_3\,, \\
    \widehat{S}_{\kappa,n}^{(m)}(\{\hat{q}_\ell\},\lambda_s) &=
    \widehat{S}_{\kappa,n}^{(m)}(\{\hat{q}_\ell\})
    + \sum_{r=1}^{2m-1-n}\frac{(n+r)!}{n!\,r!}\,
    \widehat{S}_{\kappa,n+r}^{(m)}(\{\hat{q}_\ell\})\,\ln^r\!\lambda_s\,, \\
    \mathcal{I}_{i;\ell=1, k,n}^{(m)}(z,\lambda_F,\lambda_a) &=
    \mathcal{I}_{ak,n}^{(m)}(z,\lambda_F)
    + \sum_{r=1}^{2m-1-n}\frac{(n+r)!}{n!\,r!}\,
    \mathcal{I}_{ak,n+r}^{(m)}(z,\lambda_F)\,\ln^r\!\lambda_a\,,
\end{align}
\end{subequations}
and analogously for $\mathcal{I}_{bk,n}^{(m)}(z,\lambda_F,\lambda_b)$. 
The $\lambda$
parameters are defined as
\begin{equation}
    \lambda_3 = \frac{Q_3 M}{\mu_R^2}\,, \qquad
    \lambda_s = \frac{M}{\mu_R}\,, \qquad
    \lambda_1 = \frac{Q_1 M}{\mu_R}\,, \qquad
    \lambda_2 = \frac{Q_2 M}{\mu_R}\,, \qquad
    \lambda_F = \frac{\mu_R^2}{\mu_F^2}\,.
\end{equation}


\subsection{NNLO ingredients}
\label{subsec:NNLO-ingred}

Unlike the NLO case, there is no compact master expression for the
NNLO subtraction coefficients. For brevity the scale arguments are suppressed throughout this subsection. The contributions which make up these NNLO coefficients are organised into three types of quantity. The genuine $m$-loop contributions from the jet, beam, and soft functions as
\begin{equation}
\begin{aligned}
\label{eq:Xn-def}
    X_n^{(m)} \equiv &J_{c,n}^{(m)} + \widehat{S}_{\kappa,n}^{(m)} \\
&   + \sum_k\int\frac{dz}{z}\,\mathcal{I}_{ak,n}^{(m)}\frac{f_k(x_a/z,\mu_F)}{f_a(x_a,\mu_F)}
    + \sum_k\int\frac{dz}{z}\,\mathcal{I}_{bk,n}^{(m)}\frac{f_k(x_b/z,\mu_F)}{f_b(x_b,\mu_F)}\,.
    \end{aligned}
\end{equation}
The cross terms involving the one-loop
virtual amplitude $\vec{C}^{(1)}$ are
\begin{equation}
\label{eq:Xn11-def}
\begin{aligned}
    X_n^{(1+1)} \equiv& \,H^{(1)}_{\rm SCET}
    \left[J_{c,n}^{(1)} + \widehat{S}_{\kappa,n}^{(1)} \right.\\
  &\left.  + \sum_k\int\frac{dz}{z}\,\mathcal{I}_{ak,n}^{(1)}\frac{f_k(x_a/z,\mu_F)}{f_a(x_a,\mu_F)}
    + \sum_k\int\frac{dz}{z}\,\mathcal{I}_{bk,n}^{(1)}\frac{f_k(x_b/z,\mu_F)}{f_b(x_b,\mu_F)}\right]\,.
    \end{aligned}
\end{equation}
Finally, the cross-terms between two one-loop functions arising from the
$\mathcal{L}_n\otimes\mathcal{L}_m$ convolution are
\begin{equation}
\label{eq:Xnm11-def}
\begin{aligned}
    X_{n,m}^{(1+1)} \equiv{}&
    \left[J_{c,n}^{(1)}\,\widehat{S}_{\kappa,m}^{(1)}\right.\\
    &+\!\left(J_{c,m}^{(1)}+\widehat{S}_{\kappa,m}^{(1)}\right)\!
    \left(\sum_k\int\frac{dz}{z}\,\mathcal{I}_{ak,n}^{(1)}
    \frac{f_k(x_a/z,\mu_F)}{f_a(x_a,\mu_F)}
    +\sum_k\int\frac{dz}{z}\,\mathcal{I}_{bk,n}^{(1)}
    \frac{f_k(x_b/z,\mu_F)}{f_b(x_b,\mu_F)}\right)\\
    &\left.+\sum_k\int\frac{dz}{z}\,\mathcal{I}_{ak,n}^{(1)}
    \frac{f_k(x_a/z,\mu_F)}{f_a(x_a,\mu_F)}
    \cdot\sum_{k'}\int\frac{dz'}{z'}\,\mathcal{I}_{bk',m}^{(1)}
    \frac{f_{k'}(x_b/z',\mu_F)}{f_b(x_b,\mu_F)}\right]\,.
\end{aligned}
\end{equation}
With these definitions, the NNLO subtraction coefficients are
\begin{subequations}
\label{eq:NNLO-Cn}
\begin{align}
    \mathcal{C}_3^{(2)} &= X_3^{(2)} + X_{1,1}^{(1+1)}\,, 
    \label{eq:C32}\\
    \mathcal{C}_2^{(2)} &= X_2^{(2)} + \frac{3}{2}\!\left(
        X_{0,1}^{(1+1)} + X_{1,0}^{(1+1)}\right)\,,
    \label{eq:C22}\\
    \mathcal{C}_1^{(2)} &= X_1^{(2)} + X_1^{(1+1)} + 2X_{0,0}^{(1+1)}
        - \frac{\pi^2}{3}X_{1,1}^{(1+1)}+X_{1,-1}^{(1+1)}+X_{-1,1}^{(1+1)}\,,
    \label{eq:C12}\\
    \mathcal{C}_0^{(2)} &= X_0^{(2)} + X_0^{(1+1)}
        - \frac{\pi^2}{6}\!\left(X_{0,1}^{(1+1)}+X_{1,0}^{(1+1)}\right)
        + 2\zeta_3\,X_{1,1}^{(1+1)}+X_{0,-1}^{(1+1)}+X_{-1,0}^{(1+1)}\,.
    \label{eq:C02}
\end{align}
\end{subequations}
The last two terms Eqs.~\eqref{eq:C12} and~\eqref{eq:C02} are not present in Ref.~\cite{Gaunt:2015pea}.

\section{SCET subtraction ingredients: Appendix}
\label{ap:N-jet-ing}
 Legs are labelled $a$, $b$ (incoming) and
$c$ (outgoing jet), in place of $\ell=1,2,3$ used
when calculating the \textsc{ARES} coefficients above.

\subsection{NLO subtraction ingredients}
\begin{subequations}
\label{eq:C11-ingred}
\begin{align}
&J_{c,1}^{(1)} = 4C_c\,, \\[1ex]
&\mathcal{I}_{ak,1}^{(1)} = 4C_a\,\delta_{ak}\,\delta(1-z)\,, \\[1ex]
&\mathcal{I}_{bk,1}^{(1)} = 4C_b\,\delta_{bk}\,\delta(1-z)\,, \\[1ex]
&\widehat{S}_{\kappa,1}^{(1)} = -8(C_a+C_b+C_c)\,.
\end{align}
\end{subequations}
%
\begin{subequations}
\label{eq:C10-ingred}
\begin{align}
&J_{c,0}^{(1)} = -2\gamma_c^{(0)}\,, \\[1ex]
&\mathcal{I}_{ak,0}^{(1)} = -2\gamma_a^{(0)}\,\delta_{ak}\,\delta(1-z)+2P_{ak}\,, \\[1ex]
&\mathcal{I}_{bk,0}^{(1)} = -2\gamma_b^{(0)}\,\delta_{bk}\,\delta(1-z)+2P_{bk}\,, \\[1ex]
&\widehat{S}_{\kappa,0}^{(1)} =
    4C_a\!\left(\ln\frac{Q_{ab}^2}{Q_a Q_b}+\ln\frac{Q_{ac}^2}{Q_a Q_c}
    -\ln\frac{Q_{bc}^2}{Q_b Q_c}\right) \nonumber\\[0.5ex]
  &\qquad+4C_b\!\left(\ln\frac{Q_{ab}^2}{Q_a Q_b}+\ln\frac{Q_{bc}^2}{Q_b Q_c}
    -\ln\frac{Q_{ac}^2}{Q_a Q_c}\right) \nonumber\\[0.5ex]
  &\qquad+4C_c\!\left(\ln\frac{Q_{ac}^2}{Q_a Q_c}+\ln\frac{Q_{bc}^2}{Q_b Q_c}
    -\ln\frac{Q_{ab}^2}{Q_a Q_b}\right)\,.
\end{align}
\end{subequations}
%
\begin{subequations}
\label{eq:Cd1-ingred}
\begin{align}
&J_{c,-1}^{(1)} = \begin{cases}
    C_F(7-\pi^2)\,, & \text{if $c$ is a quark}\,, \\[4pt]
    C_A\!\left(\dfrac{4}{3}-\pi^2\right)+\dfrac{20}{3}\pi\beta_0\,, &\text{if $c$ is a gluon}\,,
\end{cases} \\[1ex]
&\mathcal{I}_{ak,-1}^{(1)} = 2I_{ak}^{(1)}+2\ln\lambda_F\,P_{ak}\,, \\[1ex]
&\mathcal{I}_{bk,-1}^{(1)} = 2I_{bk}^{(1)}+2\ln\lambda_F\,P_{bk}\,, \\[1ex]
&\widehat{S}_{\kappa,-1}^{(1)}
  = C_a\!\left(\ln^2\frac{Q_{bc}^2}{Q_b Q_c}
     -\ln^2\frac{Q_{ab}^2}{Q_a Q_b}
     -\ln^2\frac{Q_{ac}^2}{Q_a Q_c}+\frac{\pi^2}{6}\right) \notag\\[0.5ex]
  &\qquad+C_b\!\left(\ln^2\frac{Q_{ac}^2}{Q_a Q_c}
     -\ln^2\frac{Q_{ab}^2}{Q_a Q_b}
     -\ln^2\frac{Q_{bc}^2}{Q_b Q_c}+\frac{\pi^2}{6}\right) \notag\\[0.5ex]
  &\qquad+C_c\!\left(\ln^2\frac{Q_{ab}^2}{Q_a Q_b}
     -\ln^2\frac{Q_{ac}^2}{Q_a Q_c}
     -\ln^2\frac{Q_{bc}^2}{Q_b Q_c}+\frac{\pi^2}{6}\right) \notag\\[0.5ex]
  &\qquad+2C_a\!\left(I_{bc,a}+I_{cb,a}-I_{ab,c}-I_{ba,c}-I_{ac,b}-I_{ca,b}\right) \notag\\[0.5ex]
  &\qquad+2C_b\!\left(I_{ac,b}+I_{ca,b}-I_{ab,c}-I_{ba,c}-I_{bc,a}-I_{cb,a}\right) \notag\\[0.5ex]
  &\qquad+2C_c\!\left(I_{ab,c}+I_{ba,c}-I_{ac,b}-I_{ca,b}-I_{bc,a}-I_{cb,a}\right)\,.
\end{align}
\end{subequations}
The matching functions, $I_{ak}^{(1)}$ and $I_{bk}^{(1)}$, are expressed as:
\begin{subequations}
\begin{align}
  I^{(1)}_{qq}(z) &=C_F\left[ \left(\frac{\ln(1-z)}{1-z}\right)_+(1+z^2)
  -\frac{\pi^2}{6}\delta(1-z)+\Theta(1-z)\left(1-z-\frac{1+z^2}{1-z}\ln z\right)\right]\,,\\[1ex]
  I^{(1)}_{qg}(z) &= T_R\left[\left(z^2+(1-z)^2\right)\ln\frac{1-z}{z}
  +2z(1-z)\right]\,,\\[1ex]
  I^{(1)}_{gq}(z) &= C_F\left[\frac{1+(1-z)^2}{z}\ln\frac{1-z}{z}+z\right]\,,\\[1ex]
  I^{(1)}_{gg}(z) &= C_A\left[\left(\frac{\ln(1-z)}{1-z}\right)_+\frac{2(1-z+z^2)^2}{z}-\frac{\pi^2}{6}\delta(1-z)-\frac{2\ln z}{(1-z)_+}\frac{(1-z+z^2)^2}{z}\right]\notag\\[0.5ex]
  &\qquad\qquad+\frac{\beta_0}{2}\delta(1-z)\,.
\end{align}
\end{subequations}

\subsection{NNLO subtraction ingredients}
\begin{subequations}
\label{eq:C23-ingred}
\begin{align}
&J_{c,3}^{(2)} = 8C_c^2\,, \\[1ex]
&\mathcal{I}_{ak,3}^{(2)} = 8C_a^2\,\delta_{ak}\,\delta(1-z)\,, \\[1ex]
&\mathcal{I}_{bk,3}^{(2)} = 8C_b^2\,\delta_{bk}\,\delta(1-z)\,, \\[1ex]
&\widehat{S}_{\kappa,3}^{(2)} = 32\left(C_a+C_b+C_c\right)^2\,.
\end{align}
\end{subequations}
\begin{subequations}
\label{eq:C22-ingred}
\begin{align}
&J_{c,2}^{(2)} = -2C_c\left(6\gamma_c^{(0)}+4\pi\beta_0\right)\,, \\[1ex]
&\mathcal{I}_{ak,2}^{(2)} = 4C_a\left[
    -\left(3\gamma_a^{(0)}+2\pi\beta_0\right)\delta_{ak}\,\delta(1-z)
    +3P_{ak}\right]\,, \\[1ex]
&\mathcal{I}_{bk,2}^{(2)} = 4C_b\left[
    -\left(3\gamma_b^{(0)}+2\pi\beta_0\right)\delta_{bk}\,\delta(1-z)
    +3P_{bk}\right]\,, \\[1ex]
&\widehat{S}_{\kappa,2}^{(2)} =
    4(C_a+C_b+C_c)\left[8\pi\beta_0-12C_a\!\left(\ln\frac{Q_{ab}^2}{Q_a Q_b}
    +\ln\frac{Q_{ac}^2}{Q_a Q_c}-\ln\frac{Q_{bc}^2}{Q_b Q_c}\right)\right. \notag\\[0.5ex]
  &\qquad\qquad\qquad\qquad\qquad-12C_b\!\left(\ln\frac{Q_{ab}^2}{Q_a Q_b}\notag
      +\ln\frac{Q_{bc}^2}{Q_b Q_c}-\ln\frac{Q_{ac}^2}{Q_a Q_c}\right) \notag\\[0.5ex]
  &\left.\qquad\qquad\qquad\qquad\qquad-12C_c\!\left(\ln\frac{Q_{ac}^2}{Q_a Q_c}
    +\ln\frac{Q_{bc}^2}{Q_b Q_c}-\ln\frac{Q_{ab}^2}{Q_a Q_b}\right)\right]\,.
\end{align}
\end{subequations}

\begin{subequations}
\label{eq:C21-ingred}
\begin{align}
&J_{c,1}^{(2)} = 8C_c K^{(1)}-\frac{8\pi^2}{3}C_c^2
    +2\gamma_c^{(0)}\!\left(2\gamma_c^{(0)}+4\pi\beta_0\right)
    +4C_c J_{c,-1}^{(1)}\,, \\[1ex]
&\mathcal{I}_{ak,1}^{(2)} =
    \left[8C_a K^{(1)}-\frac{8\pi^2}{3}C_a^2
    +2\gamma_a^{(0)}\left(2\gamma_a^{(0)}+4\pi\beta_0\right)\right]
    \delta_{ak}\,\delta(1-z) \notag\\[0.5ex]
  &\qquad+8C_a\left(I_{ak}^{(1)}+\ln\lambda_F\,P_{ak}\right)
    +4\sum_m P_{am}\otimes_z P_{mk}
    -2\left(4\gamma_a^{(0)}+4\pi\beta_0\right)P_{ak}\,, \\[1ex]
&\mathcal{I}_{bk,1}^{(2)} =
    \left[8C_b K^{(1)}-\frac{8\pi^2}{3}C_b^2
    +2\gamma_b^{(0)}\left(2\gamma_b^{(0)}+4\pi\beta_0\right)\right]
    \delta_{bk}\,\delta(1-z) \notag\\[0.5ex]
  &\qquad+8C_b\left(I_{bk}^{(1)}+\ln\lambda_F\,P_{bk}\right)
    +4\sum_m P_{bm}\otimes_z P_{mk}
    -2\left(4\gamma_b^{(0)}+4\pi\beta_0\right)P_{bk}\,, \\[1ex]
&\widehat{S}_{\kappa,1}^{(2)} =-\left(C_a+C_b+C_c\right)\left(8\widehat{S}_{\kappa,-1}^{(1)}
    +16K^{(1)}+\frac{32\pi^2}{3}\left(C_a+C_b+C_c\right)\right) \notag \\[0.5ex]
    &\qquad+16\left[C_a\!\left(\ln\frac{Q_{ab}^2}{Q_a Q_b}
    +\ln\frac{Q_{ac}^2}{Q_a Q_c}-\ln\frac{Q_{bc}^2}{Q_b Q_c}\right)\right. \notag\\[0.5ex]
  &\qquad\qquad+C_b\!\left(\ln\frac{Q_{ab}^2}{Q_a Q_b}
    +\ln\frac{Q_{bc}^2}{Q_b Q_c}-\ln\frac{Q_{ac}^2}{Q_a Q_c}\right) \notag\\[0.5ex]
  &\left.\qquad\qquad+C_c\left(\ln\frac{Q_{ac}^2}{Q_a Q_c}
    +\ln\frac{Q_{bc}^2}{Q_b Q_c}
    -\ln\frac{Q_{ab}^2}{Q_a Q_b}\right)\right]^2\notag\\[0.5ex]
  &\qquad+8\pi\beta_0\left[-4C_a\!\left(\ln\frac{Q_{ab}^2}{Q_a Q_b}
    +\ln\frac{Q_{ac}^2}{Q_a Q_c}-\ln\frac{Q_{bc}^2}{Q_b Q_c}\right)\right. \notag\\[0.5ex]
  &\qquad\qquad\qquad-4C_b\!\left(\ln\frac{Q_{ab}^2}{Q_a Q_b}
    +\ln\frac{Q_{bc}^2}{Q_b Q_c}-\ln\frac{Q_{ac}^2}{Q_a Q_c}\right) \notag\\[0.5ex]
  &\left.\qquad\qquad\qquad-4C_c\!\left(\ln\frac{Q_{ac}^2}{Q_a Q_c}
    +\ln\frac{Q_{bc}^2}{Q_b Q_c}
    -\ln\frac{Q_{ab}^2}{Q_a Q_b}\right)\right]\,.
\end{align}
\end{subequations}

\begin{subequations}
\label{eq:C20-ingred}
\begin{align}
&J_{c,0}^{(2)} = 16C_c^2\zeta_3+\frac{4\pi^2}{3}C_c\gamma_c^{(0)}
    -\left(4\gamma_c^{(0)}+4\pi\beta_0\right)J_{c,-1}^{(1)}
    -\frac{\gamma_{c,s}^{(1)}}{2}\,, \\[1ex]
&\mathcal{I}_{ak,0}^{(2)} =
    \left(16C_a^2\zeta_3+\frac{4\pi^2}{3}C_a\gamma_a^{(0)}
    -\frac{\gamma_{a,s}^{(1)}}{2}\right)\delta_{ak}\,\delta(1-z)
    -\frac{4\pi^2}{3}C_a P_{ak} \notag\\[0.5ex]
  &\qquad-\left(4\gamma_a^{(0)}+8\pi\beta_0\right)I_{ak}^{(1)}
    +4\sum_m I_{am}^{(1)}\otimes_z P_{mk}+4P_{ak}^{(1)} \notag\\[0.5ex]
  &\qquad+\ln\lambda_F\!\left[-4\gamma_a^{(0)}P_{ak}
    +4\sum_m P_{am}\otimes_z P_{mk}+4P_{ak}^{(1)}\right]\,, \\[1ex]
&\mathcal{I}_{bk,0}^{(2)} =
    \left(16C_b^2\zeta_3+\frac{4\pi^2}{3}C_b\gamma_b^{(0)}
    -\frac{\gamma_{b,s}^{(1)}}{2}\right)\delta_{bk}\,\delta(1-z)
    -\frac{4\pi^2}{3}C_b P_{bk} \notag\\[0.5ex]
  &\qquad-\left(4\gamma_b^{(0)}+8\pi\beta_0\right)I_{bk}^{(1)}
    +4\sum_m I_{bm}^{(1)}\otimes_z P_{mk}+4P_{bk}^{(1)} \notag\\[0.5ex]
  &\qquad+\ln\lambda_F\!\left[-4\gamma_b^{(0)}P_{bk}
    +4\sum_m P_{bm}\otimes_z P_{mk}+4P_{bk}^{(1)}\right]\,, \\[1ex]
&\widehat{S}_{\kappa,0}^{(2)} =
    64\left(C_a+C_b+C_c\right)^2\zeta_3-\gamma_S^{(1)}\left(C_a+C_b+C_c\right)
    -8\pi\beta_0\,\widehat{S}_{\kappa,-1}^{(1)} \notag\\[0.5ex]
  &\qquad+\frac{4\pi^2}{3}(C_a+C_b+C_c)\left[
    4C_a\!\left(\ln\frac{Q_{ab}^2}{Q_a Q_b}
    +\ln\frac{Q_{ac}^2}{Q_a Q_c}-\ln\frac{Q_{bc}^2}{Q_b Q_c}\right)\right. \notag\\[0.5ex]
  &\qquad\qquad\qquad\qquad\qquad\qquad+4C_b\!\left(\ln\frac{Q_{ab}^2}{Q_a Q_b}
    +\ln\frac{Q_{bc}^2}{Q_b Q_c}-\ln\frac{Q_{ac}^2}{Q_a Q_c}\right) \notag\\[0.5ex]
  &\left.\qquad\qquad\qquad\qquad\qquad\qquad+4C_c\!\left(\ln\frac{Q_{ac}^2}{Q_a Q_c}
    +\ln\frac{Q_{bc}^2}{Q_b Q_c}
    -\ln\frac{Q_{ab}^2}{Q_a Q_b}\right)\right] \notag\\[0.5ex]
  &\qquad+\left(\widehat{S}_{\kappa,-1}^{(1)}+K^{(1)}\right)\!\left[
    4C_a\!\left(\ln\frac{Q_{ab}^2}{Q_a Q_b}
    +\ln\frac{Q_{ac}^2}{Q_a Q_c}-\ln\frac{Q_{bc}^2}{Q_b Q_c}\right)\right. \notag\\[0.5ex]
  &\qquad\qquad\qquad\qquad\qquad+4C_b\!\left(\ln\frac{Q_{ab}^2}{Q_a Q_b}
    +\ln\frac{Q_{bc}^2}{Q_b Q_c}-\ln\frac{Q_{ac}^2}{Q_a Q_c}\right) \notag\\[0.5ex]
  &\left.\qquad\qquad\qquad\qquad\qquad+4C_c\!\left(\ln\frac{Q_{ac}^2}{Q_a Q_c}
    +\ln\frac{Q_{bc}^2}{Q_b Q_c}
    -\ln\frac{Q_{ab}^2}{Q_a Q_b}\right)\right]\,.
\end{align}
\end{subequations}
The convolution $\otimes_z$ is defined as:
\begin{equation}
P_{ij}(z)\otimes_z P_{jk}(z)=\int_z^1\frac{dw}{w}P_{ij}(w)P_{jk}\left(\frac{z}{w}\right)
\end{equation}
and analogously for $I_{ij}^{(1)}\otimes_z P_{jk}$, where the index $j$ is not summed over. 
Additionally,  the noncusp soft anomalous dimension is
\begin{equation}
\gamma_S^{(1)}=C_A\left(-\frac{64}{9}+28\zeta_3\right)+4\pi\beta_0\left(-\frac{56}{9}+\frac{\pi^2}{3}\right)\,,
\end{equation}
while those for the beam and jet functions $\gamma^{(1)}_{\ell,s}$ are, for a quark,
\begin{equation}
\gamma_{q,s}^{(1)}=C_F\left[C_A\left(\frac{146}{9}-80\zeta_3\right)+C_F(3-4\pi^2+48\zeta_3)+4\pi\beta_0\left(\frac{121}{9}+\frac{2\pi^2}{3}\right)\right]\,,
\end{equation}
and for a gluon,
\begin{equation}
\gamma_{g,s}^{(1)}=C_A\left[C_A\left(\frac{182}{9}-32\zeta_3\right)+4\pi\beta_0\left(\frac{94}{9}-\frac{2\pi^2}{3}\right)\right]\,.
\end{equation}
The subscript $s$ indicates that these follow the $N$-jettiness SCET convention opposed to 
those used within the \textsc{ARES} formalism as defined in Eq.~\eqref{eq:gammal1-flav}.


\section{ISR parts}
\subsection{qg channel}
The contribution of the splitting $g\to q\bar q$ can be obtained in a
similar way, with the major difference that there is no soft limit to
subtract and there are no virtual corrections.
This gives three contributions, $C_{\mathrm{hc},qg}^{(1)}$, $\delta\mathcal{F}_{\mathrm{hc},qg}$, $\delta\mathcal{F}_{\mathrm{rec},qg}$:
  \begin{multline}
    \label{eq:C1qg-final}
C^{(1)}_{\mathrm{hc}, qg}=  \int_{x_\ell}^1\frac{dz}{z}\,\frac{f_g(x_\ell/z,Qv^{\frac{1}{a+b_\ell}})}{f_q(x_\ell,Qv^{\frac{1}{a+b_\ell}})}\times \\ \times
\left[\frac{2 }{a+b_\ell}\left(b_\ell\left(\frac{\ln(1-z)}{1-z}\right)_+-\frac{1}{(1-z)_+}
        \left(\langle\ln d_\ell g_\ell\rangle-b_\ell \ln\frac{2E_\ell}{Q}\right)\right)(1-z)\hat P_{qg}(z;0)-P_{qg}^{(\varepsilon)}(z)\right]\,.
  \end{multline}
\begin{equation}
\label{eq:dFrec-qg}
\begin{aligned}
\delta \mathcal{F}_{\mathrm{rec},qg} = \frac{\alpha_s\left(Q v^{1/(a+b_\ell)}\right)}{\alpha_s(Q)\left(a+b_\ell\right)}\int_0^\infty\frac{d\zeta}{\zeta} &\int_0^{2\pi}\frac{d\phi}{2\pi}\int \dZ\int_{x_\ell}^1 \frac{dz}{z} \hat P_{qg}(z;0)\,\frac{f_g(x_\ell/z,Qv^{1/(a+b_\ell)})}{f_q(x_\ell,Qv^{1/(a+b_\ell)})} \\
&\times  
\left[\Theta\left(1-\frac{V_{\mathrm{hc}}(\{\tilde{p}\},k,\{k_i\})}{v}\right)-\Theta\left(1-\frac{V_{\mathrm{sc}}(\{\tilde{p}\},k,\{k_i\})}{v}\right)\right]\,,
\end{aligned}
\end{equation}
\begin{multline}
\label{eq:dFhc-qg}
\delta \mathcal{F}_{\mathrm{hc},qg} = \frac{\alpha_s\left(Qv^{1/(a+b_\ell)}\right)}{\alpha_s(Q)\left(a+b_\ell\right)}\left(\int_{x_\ell}^1 \frac{dz}{z}\frac{f_g(x_\ell/z,Qv^{1/(a+b_\ell)})}{f_q(x_\ell,Qv^{1/(a+b_\ell)})} P_{qg}(z)\right)\int_0^\infty\frac{d\zeta}{\zeta}\int_0^{2\pi}\frac{d\phi}{2\pi}\times \\
\times \int \dZ   \left[\Theta\left(1-\frac{V_{\mathrm{sc}}(\{\tilde{p}\},k,\{k_i\})}{v}\right) -\Theta(1-\zeta)\Theta\left(1-\frac{V_{\mathrm{sc}}(\{\tilde{p}\},\{k_i\})}{v}\right)\right]\,.
\end{multline}
Where we have introduced the splitting functions
\begin{equation}
  \label{eq:Pqg}
  \hat P_{qg}(z;\varepsilon) = \hat P_{qg}(z;0)+\varepsilon \hat P_{qg}^{(\varepsilon)}(z)\,,\quad
 \hat P_{qg}(z;0)= P_{qg}(z) = T_R \left[z^2+(1-z)^2\right]\,,
   \quad \hat P_{qg}^{(\varepsilon)}(z) = -2T_R z(1-z)\,.
 \end{equation}


\section{Hard function}

\chapter{One-loop hard function}
\label{ap:H1-extraction}

Following the argument of Sec.~\ref{subsec:NNLL-coeff}, we extract the
hard coefficient $H^{(1)}_{ij}$ of Eq.~\eqref{eq:HQ-expanded} and the
one-loop hard function $H^{(1)}_{\rm SCET}$ of Eq.~\eqref{eq:CH-SCET}
for the $q\bar q\to Hg$ channel. The remaining partonic channels follow
analogously. Since the relation between the two is fixed by the colour
charges and directions of the external partons alone
(cf.\ Eq.~\eqref{eq:HQ}), the finite remainder obtained below is
specific to Higgs production, but the relation it validates is not.

\paragraph{Pole comparison.} Expanding Eq.~\eqref{eq:HQ-expanded} in $\varepsilon$ for the
$q\bar{q}\to Zg$ process, the double-pole
coefficient is
\begin{equation}
  \label{eq:H-pole2}
  -\frac{1}{\varepsilon^2}(C_A+2C_F)\,,
\end{equation}
and the single-pole coefficient is
\begin{equation}
  \label{eq:H-pole1}
  \frac{1}{\varepsilon}\!\left(
  -\frac{11}{6}C_A-3C_F+\frac{n_f}{3}
  +(C_A-2C_F)\ln\frac{\mu_R^2}{Q_{12}^2}
  -C_A\ln\frac{\mu_R^2}{Q_{13}^2}
  -C_A\ln\frac{\mu_R^2}{Q_{23}^2}
  \right)\,.
\end{equation}
Comparing with the poles found in Ref.~\cite{Ravindran:2002dc}. The bare differential cross section at leading-order is
\begin{equation}
\label{eq:LO-bare}
  Q_{12}^4\frac{d^2\sigma_{0,\,q\bar{q}\to g H}}{dQ_{13}^2\,dQ_{23}^2} =
  \frac{\pi\,\delta(Q_{12}^2-Q_{13}^2-Q_{23}^2-m_H^2)}{\Gamma(1-\varepsilon)}
  \,\frac{1}{4N_c^2}
  \left(\frac{\mu_0^2\,Q_{12}^2}{Q_{13}^2Q_{23}^2}\right)^\varepsilon
 \left( \frac{\alpha_s(\mu_0^2)}{4\pi}\right)^3
 |\mathcal{M}^{(0)}_{q\bar{q}\to g H}|^2\,,
\end{equation}
where $\mu_0$ is the bare reference scale.

We must renormalise the strong coupling constant due to UV divergences using,
\begin{equation}
  \label{eq:Catani-renorm}
  \alpha_s(\mu_0^2) =
\alpha_s(\mu_R^2)
  \!\left[1-\frac{\alpha_s(\mu_R^2)\beta_0}{\varepsilon}
  +\mathcal{O}(\alpha_s^2)\right] \left(\frac{\mu_R^2}{\mu_0^2}\right)^\varepsilon\,,
\end{equation}
leading to the renormalised form of the LO differential cross-section
\begin{equation}
\begin{aligned}
  \label{eq:LO-ren}
  Q_{12}^4\frac{d^2\sigma_{q\bar{q}\to gH}^{{\rm LO},\rm Ren}}{dQ_{13}^2\,dQ_{23}^2}& \simeq
  \pi\delta(Q_{12}^2-Q_{13}^2-Q_{23}^2-M^2)\,
  \frac{S_\varepsilon}{\Gamma(1-\varepsilon)}
  \,\frac{1}{4N_c^2}
  \left(\frac{\mu_R^2\,Q_{12}^2}{Q_{13}^2Q_{23}^2}\right)^\varepsilon
\left(\frac{\alpha_s(\mu_R^2)}{4\pi}\right)^3
\left(\frac{\mu_R^2}{\mu_0^2}\right)^{2\varepsilon}\\
&\times  \!\left[1-\frac{3\alpha_s(\mu_R^2)\beta_0}{\varepsilon}\right] |M_{q\bar{q}\to gH}^{(1)}|^2\,.
  \end{aligned}
  \end{equation}
 
The bare one-loop virtual correction have the form
\begin{equation}
\label{eq:VIRT-bare}
\begin{aligned}
  &Q_{12}^4\frac{d^2\hat\sigma_{q\bar{q}\to g H}^{\rm VIRT}}{dQ_{13}^2\,dQ_{23}^2} =
  \pi\delta(Q_{12}^2-Q_{13}^2-Q_{23}^2-M^2)\,S_\varepsilon^2\,\frac{1}{N_c^2}
 \left(\frac{\alpha_s(\mu_0^2)}{4\pi}\right)^2
   \left(\frac{\alpha_s(\mu_R^2)}{4\pi}\right)^2  \\[0.5ex]
  &\quad\qquad\qquad\qquad\times\left[ n_f\!\left(-\frac{1}{3\varepsilon}
  +\frac{1}{3}\ln\frac{Q_{23}^2}{\mu_R^2}
  +\frac{1}{3}\ln\frac{Q_{13}^2}{\mu_R^2}
  -\frac{5}{9}\right) \right.\\
   &\qquad\qquad\qquad\qquad+C_A\!\left(-\frac{1}{2\varepsilon^2}
  +\!\left(-\ln\frac{Q_{12}^2}{\mu_R^2}
  +\ln\frac{Q_{23}^2}{\mu_R^2}
  +\ln\frac{Q_{13}^2}{\mu_R^2}
  +\frac{11}{6}\right)\frac{1}{\varepsilon}\right.\\[0.5ex]
  &\qquad\qquad\qquad\qquad+\operatorname{Li}_2\!\left(
  \frac{Q_{12}^2-M^2}{Q_{12}^2}\right)
  -\ln\frac{Q_{23}^2}{\mu_R^2}\ln\frac{Q_{13}^2}{\mu_R^2}
  +\ln\frac{Q_{23}^2}{\mu_R^2}\ln\frac{Q_{12}^2}{\mu_R^2}
  +\ln\frac{Q_{13}^2}{\mu_R^2}\ln\frac{Q_{12}^2}{\mu_R^2}\\[0.5ex]
  &\qquad\qquad\qquad\qquad\left.-\ln^2\frac{Q_{23}^2}{\mu_R^2}
  -\ln^2\frac{Q_{13}^2}{\mu_R^2}
  -\frac{11}{6}\ln\frac{Q_{23}^2}{\mu_R^2}
  -\frac{11}{6}\ln\frac{Q_{13}^2}{\mu_R^2}
  -\frac{7}{2}\zeta_2+\frac{38}{9}\right)\\[0.5ex]
  &\qquad\qquad\qquad\qquad+C_F\left(-\frac{1}{\varepsilon^2}
  +\!\left(\ln\frac{Q_{23}^2}{\mu_R^2}
  +\ln\frac{Q_{13}^2}{\mu_R^2}-\frac{3}{2}\right)\frac{1}{\varepsilon}
  +\operatorname{Li}_2\!\left(\frac{-Q_{23}^2}{M^2}\right)
  +\operatorname{Li}_2\!\left(\frac{-Q_{13}^2}{M^2}\right)\right.\\[0.5ex]
  &\qquad\qquad\qquad\qquad-2\ln\frac{Q_{23}^2}{\mu_R^2}\ln\frac{Q_{13}^2}{\mu_R^2}
  -\frac{1}{2}\ln^2\!\left(\frac{M^2+Q_{23}^2}{Q_{23}^2}\right)
  -\frac{1}{2}\ln^2\!\left(\frac{M^2+Q_{13}^2}{Q_{13}^2}\right)\\[0.5ex]
  &\qquad\qquad\qquad\qquad\left.+\frac{1}{2}\ln^2\!\left(\frac{M^2+Q_{23}^2}{M^2}\right)
  +\frac{1}{2}\ln^2\!\left(\frac{M^2+Q_{13}^2}{M^2}\right)
  +\frac{3}{2}\ln\frac{Q_{23}^2}{\mu_R^2}
  +\frac{3}{2}\ln\frac{Q_{13}^2}{\mu_R^2}
  +9\zeta_2-4\right)\\[0.5ex]
  &\qquad\qquad\qquad\qquad\left.+\frac{3}{2}\ln\frac{Q_{23}^2}{\mu_R^2}
  +\frac{3}{2}\ln\frac{Q_{13}^2}{\mu_R^2}
  +9\zeta_2-4\right]|M^{(1)}_{q\bar{q}\to g H}|^2\,.
\end{aligned}
\end{equation}
We renormalise the bare coupling, again using Eq.~\eqref{eq:Catani-renorm} and find
\begin{equation}
\label{eq:virt-ren}
  Q_{12}^4\frac{d^2\hat\sigma_{q\bar{q}\to g H}^{\rm VIRT, Ren}}{dQ_{13}^2\,dQ_{23}^2} \simeq
  \pi\delta(Q_{12}^2-Q_{13}^2-Q_{23}^2-M^2)\,S_\varepsilon^2
  \,\frac{1}{N_c^2}
  \left(\frac{\alpha_s(\mu_R^2)}{4\pi}\right)^4  \left(\frac{\mu_R^2}{\mu_0^2}\right)^{2\varepsilon}\left[\ldots\right]|M^{(1)}_{q\bar{q}\to g H}|^2 \,,
\end{equation}
where the $\ldots$ represent terms found in the square brackets of Eq.~\eqref{eq:VIRT-bare}.
The relevant combination for extracting the hard function is the
sum of the renormalised LO and virtual terms at order $\mathcal{O}(\alpha_s^4)$
\ABquestion{Do I need to explain why we go up to this order of $\alpha_s$ and also why it is this combination which gives the the finite virtual corrections?}
\begin{equation}
\begin{aligned}
\label{eq:LO+virt-ren}
  &Q_{12}^4\frac{d^2\sigma_{q\bar{q}\to g H}^{(1),\rm Ren}}{dQ_{13}^2\,dQ_{23}^2}
  +Q_{12}^4\frac{d^2\hat\sigma_{q\bar{q}\to g H}^{\rm VIRT,Ren}}{dQ_{13}^2\,dQ_{23}^2}\\[0.5ex]
 &\simeq  \pi\delta(Q_{12}^2-Q_{13}^2-Q_{23}^2-M^2)\,
  \frac{S_\varepsilon}{\Gamma(1-\varepsilon)}
  \,\frac{1}{4N_c^2}
  \left(\frac{\mu_R^2\,s}{tu}\right)^\varepsilon
\left(\frac{\alpha_s(\mu_R^2)}{4\pi}\right)^4
\left(\frac{\mu_R^2}{\mu_0^2}\right)^{2\varepsilon}\\[0.5ex]
&\times\left\{-\frac{12 \pi \beta_0}{\varepsilon}+\right.+4 S_{\varepsilon}\Gamma(1-\varepsilon)
  \left(\frac{Q_{13}^2 Q_{23}^2}{\mu_R^2\,Q_{12}^2}\right)^\varepsilon \left[ n_f\!\left(-\frac{1}{3\varepsilon}
  +\frac{1}{3}\ln\frac{Q_{23}^2}{\mu_R^2}
  +\frac{1}{3}\ln\frac{Q_{13}^2}{\mu_R^2}
  -\frac{5}{9}\right) \right.\\[0.5ex]
   &\quad+C_A\!\left(-\frac{1}{2\varepsilon^2}
  +\!\left(-\ln\frac{Q_{12}^2}{\mu_R^2}
  +\ln\frac{Q_{23}^2}{\mu_R^2}
  +\ln\frac{Q_{13}^2}{\mu_R^2}
  +\frac{11}{6}\right)\frac{1}{\varepsilon}\right.\\[0.5ex]
  &\quad+\operatorname{Li}_2\!\left(
  \frac{Q_{12}^2-M^2}{Q_{12}^2}\right)
  -\ln\frac{Q_{23}^2}{\mu_R^2}\ln\frac{Q_{13}^2}{\mu_R^2}
  +\ln\frac{Q_{23}^2}{\mu_R^2}\ln\frac{Q_{12}^2}{\mu_R^2}
  +\ln\frac{Q_{13}^2}{\mu_R^2}\ln\frac{Q_{12}^2}{\mu_R^2}\\[0.5ex]
  &\quad\left.-\ln^2\frac{Q_{23}^2}{\mu_R^2}
  -\ln^2\frac{Q_{13}^2}{\mu_R^2}
  -\frac{11}{6}\ln\frac{Q_{23}^2}{\mu_R^2}
  -\frac{11}{6}\ln\frac{Q_{13}^2}{\mu_R^2}
  -\frac{7}{2}\zeta_2+\frac{38}{9}\right)\\[0.5ex]
  &\quad+C_F\left(-\frac{1}{\varepsilon^2}
  +\!\left(\ln\frac{Q_{23}^2}{\mu_R^2}
  +\ln\frac{Q_{13}^2}{\mu_R^2}-\frac{3}{2}\right)\frac{1}{\varepsilon}
  +\operatorname{Li}_2\!\left(\frac{-Q_{23}^2}{M^2}\right)
  +\operatorname{Li}_2\!\left(\frac{-Q_{13}^2}{M^2}\right)\right.\\[0.5ex]
  &\quad-2\ln\frac{Q_{23}^2}{\mu_R^2}\ln\frac{Q_{13}^2}{\mu_R^2}
  -\frac{1}{2}\ln^2\!\left(\frac{M^2+Q_{23}^2}{Q_{23}^2}\right)
  -\frac{1}{2}\ln^2\!\left(\frac{M^2+Q_{13}^2}{Q_{13}^2}\right)\\[0.5ex]
  &\quad\left.+\frac{1}{2}\ln^2\!\left(\frac{M^2+Q_{23}^2}{M^2}\right)
  +\frac{1}{2}\ln^2\!\left(\frac{M^2+Q_{13}^2}{M^2}\right)
  +\frac{3}{2}\ln\frac{Q_{23}^2}{\mu_R^2}
  +\frac{3}{2}\ln\frac{Q_{13}^2}{\mu_R^2}\right)\\[0.5ex]
  &\quad\left.\left.+9\zeta_2-4+\frac{3}{2}\ln\frac{Q_{23}^2}{\mu_R^2}
  +\frac{3}{2}\ln\frac{Q_{13}^2}{\mu_R^2}
  +9\zeta_2-4\right]\right\}|M^{(1)}_{q\bar{q}\to g H}|^2\,.
\end{aligned}
\end{equation}
\ABquestion{The $S_{\varepsilon}\Gamma(1-\varepsilon)$
term gives a $\ln 4\pi$ which is absorbed into the coupling in MS-bar scheme which
is what the Catani poles are in so in essence this can also be absorbed into
the coupling?}

The double-pole coefficients in Eq.~\eqref{eq:LO+virt-ren} are:
\begin{equation}
  \label{eq:vN-pole2}
  -\frac{2}{\varepsilon^2}(C_A+2C_F)\,,
\end{equation}
and the single-pole coefficient is
\begin{equation}
   \label{eq:vN-pole1}
  \frac{1}{\varepsilon}\!\left(
  -\frac{11}{3}C_A-6C_F+\frac{2n_f}{3}
  +2(C_A-2C_F)\ln\frac{\mu_R^2}{Q_{12}^2}
  -2C_A\ln\frac{\mu_R^2}{Q_{13}^2}
  -2C_A\ln\frac{\mu_R^2}{Q_{23}^2}
  \right)\,.
\end{equation}

Comparing with Eqs.~\eqref{eq:H-pole2} and~\eqref{eq:H-pole1},
the poles of Ref.~\cite{Ravindran:2002dc} are exactly twice those
of $\mathcal{V}(Q;\varepsilon)$, this factor is taken into account when 
subtracting the finite terms.

\paragraph{Extraction of the finite remainder}
The finite remainder is extracted from the
$\mathcal{O}(\varepsilon^0)$ terms of Eq.~\eqref{eq:VIRT-bare}
after subtracting the predicted poles giving rise to
the SCET hard function
\begin{equation}
\label{eq:H1-SCET}
\begin{aligned}
& H^{(1)}_{\rm SCET}=n_f\left(-\frac{10}{9}-\frac{2}{3}\ln\frac{\mu_R^2}{Q_{12}^2}\right)+\pi^2\left(
3C_F-\frac{7}{6}C_A
\right)\\[0.5ex]
& \qquad+C_A\left[
\frac{76}{9}
+2\operatorname{Li}_2\left(\frac{Q_{12}^2-M^2}{Q_{12}^2}\right)
-\frac{11}{3}\ln\frac{Q_{12}^2}{\mu_R^2}
+\frac{1}{2}\ln^2\frac{\mu_R^2}{Q_{12}^2}
-\frac{1}{2}\ln^2\frac{\mu_R^2}{Q_{13}^2}
-\frac{1}{2}\ln^2\frac{\mu_R^2}{Q_{23}^2}\right.\\[0.5ex]
&\quad\qquad\qquad\left.+\ln\frac{Q_{12}^2}{Q_{13}^2}
 \ln\frac{Q_{12}^2}{Q_{23}^2}
\right]\\[0.5ex]
&\qquad+C_F\left[-8
+2\operatorname{Li}_2\left(\frac{-Q_{13}^2}{M^2}\right)
+2\operatorname{Li}_2\left(\frac{-Q_{23}^2}{M^2}\right)
-\ln^2\frac{-Q_{13}^2-M^2}{-Q_{13}^2}
-\ln^2\frac{-Q_{23}^2-M^2}{-Q_{23}^2}\right.\\
&\quad\qquad\qquad\left.
+\ln^2\frac{M^2+Q_{13}^2}{M^2}
+\ln^2\frac{M^2+Q_{23}^2}{M^2}
-\ln^2\frac{Q_{12}^2}{\mu_R^2}
+3\ln\frac{Q_{12}^2}{\mu_R^2}
+\ln^2\frac{Q_{13}^2}{Q_{23}^2}
\right]\,.
\end{aligned}
\end{equation}
For the \textsc{ARES} formalism, the hard coefficient is given by
Eq.~\eqref{eq:H1-final}, such that we obtain: 
\begin{equation}
H_{q\bar{q}}^{(1)}=\mathcal{H}_1
+3\left(C_F-\frac{C_A}{2}\right)\ln\frac{Q_{12}^2}{Q^2}
+\left(\frac{3C_A}{4}+\pi\beta_0\right)\ln\frac{Q_{13}^2}{Q^2}
+\left(\frac{3C_A}{4}+\pi\beta_0\right)\ln\frac{Q_{23}^2}{Q^2}\,.
\end{equation}
Unlike $H_{\rm SCET}^{(1)}$, which is identified directly with the finite
remainder of the one-loop amplitude, the \textsc{ARES} hard coefficient
$\mathcal H_1$ is obtained by subtracting the remaining finite
contributions appearing in Eq.~\eqref{eq:HQ-expanded} from the SCET hard
function of Eq.~\eqref{eq:H1-SCET}. Since these subtraction terms are
universal, this provides a validation of the relation between the SCET
hard function and the corresponding \textsc{ARES} hard coefficient. We
therefore obtain\footnote{Please note the factor of 2 difference due to
Eq.~\eqref{eq:coeff-convent}.}
\begin{equation}
\begin{aligned}
\mathcal{H}_1=&\frac{1}{36}\left\{-20n_f+C_A(152-21\pi^2)+18C_F(-8+3\pi^2)
-66C_A\ln\frac{Q_{12}^2}{\mu_R^2}\right.\\[1ex]
&\quad+6\left[-3C_F\ln^2\left(\frac{Q_{13}^2+M^2}{Q_{13}^2}\right)-3C_F\ln^2\left(\frac{Q_{23}^2+M^2}{Q_{23}^2}\right)\right.\\
&\quad+3C_F\ln^2\left(\frac{M^2+Q_{13}^2}{M^2}\right)+3C_F\ln^2\left(\frac{M^2+Q_{23}^2}{M^2}\right)
+9C_F\ln\frac{Q_{13}^2}{\mu_R^2}+2n_f\ln\frac{Q_{13}^2}{\mu_R^2}\\
&\quad+6C_A\ln\frac{Q_{12}^2}{\mu_R^2}\ln\frac{Q_{13}^2}{\mu_R^2}
-6C_A\ln^2\frac{Q_{13}^2}{\mu_R^2}-6C_A\ln^2\frac{Q_{23}^2}{\mu_R^2}\\
&\quad\left.+\left(-11C_A+9C_F+2n_f+6C_A\ln\frac{Q_{12}^2}{\mu_R^2}
-6(C_A+2C_F)\ln\frac{Q_{13}^2}{\mu_R^2}\right)\ln\frac{Q_{23}^2}{\mu_R^2}-6C_A\ln^2\frac{Q_{23}^2}{\mu_R^2}\right]\\[1ex]
&\quad+6\left[11C_A-9C_F-2n_f-6C_A\ln\frac{Q_{12}^2}{\mu_R^2}
+6(C_A+C_F)\left(\ln\frac{Q_{12}^2}{\mu_R^2}+\ln\frac{Q_{23}^2}{\mu_R^2}\right)\right]
\ln\frac{Q_{13}^2Q_{23}^2}{Q_{12}^2\mu_R^2}\\[1ex]
&\quad-9(C_A+2C_F)\ln^2\frac{Q_{13}^2Q_{23}^2}{Q_{12}^2\mu_R^2}\\[1ex]
&\quad+3\left[-(C_A+2C_F)\pi^2-18(C_A-2C_F)\ln\frac{\mu_R^2}{Q_{12}^2}
-6(C_A-2C_F)\ln^2\frac{\mu_R^2}{Q_{12}^2}\right.\\
&\quad\left.+3(3C_A+4\pi\beta_0)\ln\frac{\mu_R^2}{Q_{13}^2}
+6C_A\ln^2\frac{\mu_R^2}{Q_{13}^2}+3\ln\frac{\mu_R^2}{Q_{23}^2}\left(3C_A+4\pi\beta_0
+2C_A\ln\frac{\mu_R^2}{Q_{23}^2}\right)\right]\\[1ex]
&\quad\left.+36\left[C_A\operatorname{Li}_2\left(\frac{Q_{12}^2-M^2}{Q_{12}^2}\right)
+C_F\operatorname{Li}_2\left(-\frac{Q_{13}^2}{M^2}\right)
+C_F\operatorname{Li}_2\left(-\frac{Q_{23}^2}{M^2}\right)\right]\right\}\,.
\end{aligned}
\end{equation}

\section{General set up - chapter 5}

Eq.~\eqref{eq:dsigma-def}
requires a mapping between the full event -- made up
of the three hard partons $\{\tilde p\}\equiv\{\tilde p_1,\tilde p_2,\tilde p_3\}$ and an arbitrary number of soft and/or collinear emissions -- and
a Born event with only three hard partons $\{p_1,p_2,p_3\}$.
The difference between the approximation on the right-hand side of
Eq.~\eqref{eq:dsigma-def} and the full expression for
$ d\sigma(v)/(dq_T dy_Z d\eta_J)$ is a well-defined expansion in powers of the
strong coupling $\alpha_s$, with coefficients vanishing as a power of
$v$, potentially enhanced by logarithms of $v$. 

\section{NNLL cumulant: chapter 5}
The NNLL radiator $R_{\rm NNLL}(v)$ includes all virtual
  corrections of \emph{soft and/or collinear origin}. These are divergent,
  but their divergences cancel with those of real radiation. In
  order to cancel the singularities of virtual corrections analytically, the
  radiator also includes a subset of real corrections for suitably
  defined unresolved emissions. The procedure to define those
  emissions is not unique (cf.\ Sec.\ref{sec:NLL-resum}). 
  A proposal adequate for three emitting legs is given in
  Ref.~\cite{Arpino:2019ozn}, leading to a master formula for
  $R_{\rm NNLL}(v)$ whose explicit expression depends only on the
  single-emission parameters $a$, $b_\ell$, $d_\ell$, and
  $g_\ell(\phi)$. For completeness, we illustrate how
  the NNLL radiator extends from the NLL result in
  Sec.~\ref{sec:radiator}. 

\section{One-jettiness ARES}
\begin{equation}
\begin{aligned}
	&f_{\rm sc}^{(\ell\bar\ell)}(\eta^{(\ell\bar\ell)},\phi^{(\ell\bar\ell)})=\frac{Q_{\ell\bar\ell}^2}{MQ_{\ell}}e^{-\eta^{(\ell\bar\ell)}}\Theta\left(\eta^{(\ell\bar\ell)}\right)+\frac{Q_{\ell\bar\ell}^2}{MQ_{\bar\ell}}e^{\eta^{(\ell\bar\ell)}}\Theta\left(-\eta^{(\ell\bar\ell)}\right)\,,\\
	&f_{\rm wa}^{(\ell\bar\ell)}(\eta^{(\ell\bar\ell)},\phi^{(\ell\bar\ell)}) =
\min\left\{
\frac{Q_{\ell\bar\ell}^2}{MQ_1}\,e^{-\eta^{(\ell\bar\ell)}},\;
\frac{Q_{\ell\bar\ell}^2}{MQ_2}\,e^{+\eta^{(\ell\bar\ell)}},\;
\frac{Q_{\ell\bar\ell}}{M}\!\left(\cosh(\eta^{(\ell\bar\ell)}-\eta_J)
-\cos(\phi^{(\ell\bar\ell)}-\phi_J)\right)
\right\}\,.
\end{aligned}
\end{equation}


